In this chapter, we review the existing research regarding Industrial Control System (ICS) security, ranging from foundational vulnerability analysis to advanced detection frameworks. Following the logical structure of identifying research gaps, we categorize these works to highlight the necessity of our proposed standard-compliant DPI system.

\section{Vulnerability Analysis and Protocol Testing}
Understanding the inherent flaws of ICS protocols is the primary step in establishing defense. Pliatsios et al. \cite{ref_report25} provided a high-level survey on SCADA systems, categorizing vulnerabilities into protocol-level weaknesses and architectural flaws. Focusing on the specific communication medium, Chang et al. \cite{ref_report26} executed a series of vulnerability tests on the Modbus TCP protocol, demonstrating how attackers use legitimate function codes to disrupt operations. Complementing this, Luswata et al. \cite{ref_report27} utilized penetration testing to analyze SCADA security, revealing that internal lateral movement remains a critical threat vector. However, these methodologies are primarily active or manual testing approaches. They lack the capability for automated, continuous monitoring, which is essential for real-time OT network defense.

\section{Anomaly Detection and Complex Frameworks}
To address dynamic threats, several advanced frameworks have been proposed. ShadowPLCs \cite{ref_report28} is a novel scheme that mirrors the execution state of physical PLCs in a simulated environment to detect manipulation. While robust, its overall system is resource-intensive and requires high-fidelity models for every specific device. Alternatively, Kim and Park \cite{ref_report29} explored machine learning models for time-series anomaly detection. While machine learning is effective at detecting unknown threats, it often operates as a "black box," offering low interpretability for security analysts. Furthermore, while Digital Twin technologies \cite{ref_report30} offer high-fidelity simulation, they often lack a systematic mapping to international security management standards. These interpretability and standardization issues will be the focus of this thesis.

\section{Proposed Deep Packet Inspection and Standard Integration}
In contrast to high-overhead simulation or low-interpretability models, Deep Packet Inspection (DPI) provides an efficient and transparent solution. Suricata is an advanced DPI engine that utilizes a rule-based system to inspect network traffic at the application layer. When Suricata hooks into the network interface, it extracts specific fields from Modbus frames to match against predefined signatures. 
\noindent
Our core contribution lies in the systematic alignment of these detection rules with international standards. Instead of ad-hoc rule creation, our approach maps each signature to the Foundational Requirements (FRs) of ISA/IEC 62443-3-3 and the ICS Cyber Kill Chain. By forwarding these structured logs to Graylog, we provide a centralized SIEM platform that enhances situational awareness through standard-aligned alerts. This architecture ensures that technical detections are directly linked to high-level security compliance.

\section{Summary of Related Works}
In summary, while existing methods provide valuable insights into vulnerability testing and anomaly detection, they rarely integrate these technical detections with standardized security management frameworks and threat modeling methodologies like the ICS Cyber Kill Chain. To clearly illustrate the research gap, Table \ref{tab:comparison} compares these works with our proposed method.

\noindent
\begin{table}[htbp]
\caption{Research Focus and Capability Mapping of ICS Security Solutions.}
\centering
\scriptsize                          
\setlength{\tabcolsep}{3pt}          
\renewcommand{\arraystretch}{1.4}    

\begin{tabular}{@{} p{2.6cm} p{2.2cm} c c c c c c @{}}
\toprule
\textbf{Literature} & \textbf{Methodology} & 
\begin{tabular}[c]{@{}c@{}}\textbf{Continuous}\\\textbf{Monitor}\end{tabular} & 
\begin{tabular}[c]{@{}c@{}}\textbf{Interpret-}\\\textbf{ability}\end{tabular} & 
\begin{tabular}[c]{@{}c@{}}\textbf{System}\\\textbf{Overhead}\end{tabular} & 
\begin{tabular}[c]{@{}c@{}}\textbf{Unknown}\\\textbf{Threat}\\\textbf{Detection}\end{tabular} & 
\begin{tabular}[c]{@{}c@{}}\textbf{ISA/IEC}\\\textbf{62443 3-3}\end{tabular} & 
\begin{tabular}[c]{@{}c@{}}\textbf{ICS Kill}\\\textbf{Chain}\end{tabular} \\ 
\midrule

Pliatsios \cite{ref_report25} / \newline Luswata \cite{ref_report27} & Survey / \newline Pentest & No & White-box & Lightweight & No & \begin{tabular}[c]{@{}c@{}}Not\\Addressed\end{tabular} & \begin{tabular}[c]{@{}c@{}}Not\\Addressed\end{tabular} \\ \addlinespace[0.5em]

ShadowPLCs \cite{ref_report28} & State \newline Mirroring & Yes & White-box & \begin{tabular}[c]{@{}c@{}}Resource-\\intensive\end{tabular} & Yes & \begin{tabular}[c]{@{}c@{}}Not\\Addressed\end{tabular} & \begin{tabular}[c]{@{}c@{}}Not\\Addressed\end{tabular} \\ \addlinespace[0.5em]

Kim \& Park \cite{ref_report29} & Machine \newline Learning & Yes & Black-box & Moderate & Yes & \begin{tabular}[c]{@{}c@{}}Not\\Addressed\end{tabular} & \begin{tabular}[c]{@{}c@{}}Not\\Addressed\end{tabular} \\ \addlinespace[0.5em]

Digital Twin \cite{ref_report30} & Simulation & Yes & \begin{tabular}[c]{@{}c@{}}Simulation-\\dependent\end{tabular} & \begin{tabular}[c]{@{}c@{}}Resource-\\intensive\end{tabular} & Yes & \begin{tabular}[c]{@{}c@{}}Not\\Addressed\end{tabular} & \begin{tabular}[c]{@{}c@{}}Not\\Addressed\end{tabular} \\ 
\midrule[\heavyrulewidth] % 加粗分隔線，以突顯 Proposed Work

\textbf{Proposed Work} & \textbf{DPI+SIEM} & \textbf{Yes} & \textbf{Rule-based} & \textbf{Lightweight} & \textbf{Limited} & \textbf{Integrated} & \textbf{Mapped} \\ 
\bottomrule
\end{tabular}
\label{tab:comparison}
\end{table}

As shown in Table \ref{tab:comparison}, most current systems rarely integrate deep packet inspection with unified security frameworks and structured attack lifecycles. Our research addresses this exact gap by providing a lightweight, interpretable, and standard-compliant monitoring system. By mapping detection rules to both ISA/IEC 62443-3-3 Foundational Requirements and the ICS Cyber Kill Chain. This approach helps bridge packet-level detection, tactical threat awareness, and security compliance requirements.
